\section{Vragen voor de doelgroep}
\label{sec:interviews-doelgroep}

Vijf respondenten: J. De Smedt, M. Adams, M. Dhondt, M. Maratta, S. Landtsheer

\vspace{1em}

\textbf{1. Hoe vertrouwd ben je naar eigen mening met technologie en digitale toepassingen?}

Iedere respondent gaf aan vertrouwd te zijn met technologie en digitale toepassingen. Ofwel door hun huidige opleiding, of door eigen interesses en ervaringen (opgroeien met technologie).

\vspace{1em}

\textbf{2. Hoe frequent maak je gebruik van streamingdiensten?}

Alle respondenten maken minstens twee dagen in de week gebruik van streamingdiensten. Drie van de respondenten gaven aan (bijna) dagelijks te streamen.

\vspace{1em}

\textbf{3. Welke rol spelen streamingdiensten in je dagelijks leven?}

\textbf{3.1. Waarom gebruik je streamingdiensten?}

Iedere respondent gaf aan streamingdiensten te gebruiken voor entertainmentdoeleinden. Het kijken naar films en series werd het meest vernoemd. Daarnaast gaven de respondenten aan te streamen om muziek te luisteren, om reclame te vermijden en om op deze manier snel en makkelijk te kunnen kijken naar content zonder te hoeven downloaden.

\newpage

\textbf{3.2. Welke streamingdiensten gebruik je? Welke gebruik je het vaakst?}

Netflix werd door het merendeel van de respondenten het vaakst gebruikt. Daarnaast werden onder andere Spotify, Disney+, Streamz, VTM GO, GoPlay, VRT MAX, PlayTV, SoundCloud en Epidemic Sound genoemd.

\vspace{1em}

\textbf{3.3. Op welke toestellen stream je? Op welke toestellen het frequentst?}

Een meerderheid van de respondenten streamen voornamelijk op een computer of laptop. Daarnaast worden ook smartphones, televisies en Chromecast gebruikt.

\vspace{1em}

\textbf{4. Welke factoren beïnvloeden je keuze om een abonnement af te sluiten, te behouden of op te zeggen?}

De meest beïnvloedende factor bij de keuze voor een streamingabonnement is volgens de respondenten de prijs. Eveneens werd het beschikbare aanbod aangehaald als factor. Ten slotte kwam de mogelijkheid om een abonnement te delen ook aan bod, als ook of een betalend abonnement al dan niet advertenties bevat.

\vspace{1em}

\textbf{5. Wat versta je onder “dark patterns” of misleidende ontwerptechnieken op websites of apps?}

Drie respondenten beschreven dark patterns als technieken die gebruikers proberen te beïnvloeden of misleiden om bepaalde keuzes te maken in het voordeel van de streamingdienst. Twee respondenten waren niet vertrouwd met de term.

\vspace{1em}

\textbf{6. Heb je ooit het gevoel gehad dat een website of app van een streamingdienst je probeerde te sturen om een bepaalde keuze te maken? Leg uit.}

Iedere respondent kon een voorbeeld geven van een situatie waarin ze het gevoel hadden dat ze werden beïnvloed door een streamingdienst. Voornamelijk de gepersonaliseerde aanbevelingen op Netflix kwamen aan bod, maar ook de autoplay-functionaliteit werd frequent genoemd. Ten slotte werd ook het gebruik van advertenties vermeld, als ook gaf een respondent aan gestuurd te worden in de keuze van een abonnement. 

\vspace{1em}

\textbf{7. Wanneer merk je dat een streamingplatform je aandacht langer probeert vast te houden? Hoe ervaar je dat?}

Het merendeel van de respondenten merkte op dat functionaliteiten zoals autoplay (de knop om onmiddellijk naar de volgende aflevering te gaan) en automatische aanbevelingen hen langer laten kijken. Sommigen ervaren deze functionaliteiten als handig, terwijl anderen ze eerder storend of verslavend vinden. Één respondent gaf aan dit nog niet bewust te hebben opgemerkt.

\newpage

\textbf{8. Kan je een situatie beschrijven waarbij je in aanraking kwam met het (bijna) onverwacht inschrijven voor een abonnement, proefperiode of extra dienst?}

\textbf{8.1. Hoe merkte je dit op?}

Drie respondenten hadden ervaring met het in aanraking komen met onverwachte of ongewenste abonnementen. De respondenten ervoeren dit op Amazon Prime, Epidemic Sound en via de iCloud-opslag. Zij merkten dit op doordat er onverwacht geld werd afgehouden of doordat zij voortdurend werden aangespoord om te upgraden. Twee respondenten zijn niet verder gegaan met het ongewenste abonnement, één respondent viel in de val van een dark pattern waardoor een abonnement onverwacht werd verlengd. 

\vspace{1em}

\textbf{8.2. Hoe verliep het uitschrijven?}

Slechts één respondent had ervaring met het uitschrijven van een abonnement en omschreef dit als moeilijk. De overige respondenten hadden geen ervaring met het uitschrijven.

\vspace{1em}

\textbf{8.3. Wat maakte dit makkelijk of moeilijk?}

De enige respondent die een concrete ervaring had, gaf aan dat het uitschrijven verwarrend was doordat het leek alsof het abonnement al was stopgezet, terwijl er nog een extra, minder opvallende stap nodig was (een verborgen bevestigingsknop).

\vspace{1em}

\textbf{9. Welke elementen of tekenen zouden jou het vermoeden geven dat een platform je probeert te beïnvloeden?}

De respondenten noemden verschillende tekenen die kunnen wijzen op beïnvloeding, zoals opvallende knoppen, animaties, timers, tijdsgebonden aanbiedingen, veel aanbevelingen, advertenties en het aantrekkelijk voorstellen van betalende opties. Één respondent gaf aan zo'n elementen niet snel op te merken.

\vspace{2em}

\textit{Na het stellen van deze vragen, werd uitgelegd wat dark patterns exact zijn, gedemonstreerd door enkele voorbeelden.}

\vspace{2em}

\textbf{10. Stel dat er een tool bestaat die je helpt bij het herkennen van mogelijke dark patterns. Hoe zou zo'n tool gebruikers volgens jou het best kunnen helpen?}

Merendeel van de respondenten gaf aan dat de tool proactief moet waarschuwen wanneer een mogelijk dark pattern wordt gedetecteerd. Daarnaast werden suggesties gedaan zoals het weergeven van extra uitleg, het visueel aanduiden van beïnvloedende elementen, het aanbieden van een overzicht van abonnementen en het personaliseren van waarschuwingen (aantal meldingen, keuze van patronen waarvoor je gewaarschuwd wil worden).

\vspace{1em}

\textbf{11. Hoe zou je willen dat een tool jou informeert wanneer een mogelijk dark pattern wordt gedetecteerd? Enkele voorbeelden: in real-time of na het streamen een overzicht kunnen raadplegen, veel waarschuwingen of weinig...}

Merendeel van de respondenten gaven de voorkeur aan real-time meldingen, vooral wanneer een dark pattern financiële gevolgen kan hebben. Wel vonden ze dat de waarschuwingen niet te opdringerig mogen zijn en eerder subtiel weergegeven moeten worden. Enkele respondenten stelden voor om subtiele indicatoren of een interactief icoon te gebruiken in plaats van grote pop-ups. Ook het personaliseren van meldingen op basis van soort en frequentie werden meerdere keren genoemd.

\vspace{1em}

\textbf{12. Welke informatie zou je willen krijgen wanneer een mogelijk dark pattern gedetecteerd wordt?}

Zo goed als alle respondenten wilden weten welk dark pattern werd gedetecteerd en waarom dit mogelijk problematisch is. Daarnaast vonden zij het handig om te weten waar het dark pattern gedetecteerd werd, welke mogelijke gevolgen er zijn en enkele tips om ze te vermijden nuttig. 

\vspace{1em}

\textbf{13. Wanneer zou een waarschuwing volgens jou storend worden? Kan je een voorbeeld geven van een situatie waarin je een waarschuwing zou negeren?}

Bijna alle respondenten gaven aan dat waarschuwingen storend worden wanneer ze te vaak verschijnen, te groot zijn of een verplichte interactie vereisen (drukken om een melding te sluiten). Verschillende respondenten zouden waarschuwingen over bekende of minder belangrijke situaties, zoals cookies of autoplay, sneller negeren. De respondenten benadrukten daarom dat waarschuwingen relevant, subtiel en indien mogelijk personaliseerbaar moeten zijn.

\vspace{1em}


\section{Vragen voor experten}
\label{sec:interviews-experten}

Twee respondenten:

\begin{itemize}
    \item Systeemarchitect \textbf{G. De Bruyn}. Hij heeft vier jaar ervaring in front-end development bij KBC bank \& verzekeringen, waar gebruiksvriendelijkheid centraal staat.
    \item Product designer \textbf{S. Hostyn}. Hij heeft vier jaar ervaring binnen UI/UX design en werkt sinds een jaar bij Intigriti, waar frequent testen en interviews worden afgenomen bij gebruikers om het platform te verbeteren. 
\end{itemize}

\vspace{1em}

\textbf{1. In welke mate bent u vertrouwd met het concept dark patterns?}

Beide respondenten zijn vertrouwd met het concept dark patterns. Echter verschillen ze in expertise. De product designer komt er dagelijks mee in contact en beschikt over praktijkervaring, terwijl de systeemarchitect het concept voornamelijk in grote lijnen kent. Daarnaast benadrukt de product designer het spanningsveld tussen overtuigen en manipuleren, waarbij hij aangeeft dat onder andere conversiecijfers en KPI's de verleiding kunnen creëren om manipulatieve ontwerpkeuzes te maken. Volgens hem verdient het overtuigen van gebruikers met duidelijke argumenten steeds de voorkeur boven manipulatie.

\vspace{1em}

\textbf{2. Hoe kijkt u naar het probleem van dark patterns binnen streamingdiensten? Is het iets om zorgen over te maken? Waarom?}

Zowel de systeemarchitect als de product designer beschouwen dark patterns binnen streamingdiensten als een reëel probleem. Zo schuilt het gevaar erin dat de schadelijke gevolgen van de patronen vaak traag en onzichtbaar zijn en dat de manipulatieve ontwerptechnieken net onschuldig lijken. Enerzijds is er een mogelijke emotionele impact, waarbij de manipulatieve ontwerpkeuzes verslavend kijkgedrag kunnen stimuleren. Anderzijds bestaat het risico dat het gebruik van dergelijke patronen genormaliseerd wordt, waardoor de focus van streamingdiensten verschuift van het overtuigen van gebruikers naar het zo effectief mogelijk manipuleren ervan. Daarnaast benadrukken beide respondenten dat dit kan leiden tot aanzienlijke financiële gevolgen, bijvoorbeeld door ongewenste abonnementen of automatische verlengingen.

\vspace{1em}

\textbf{3. Zijn er na het lezen van de huidige draft van de bachelorproef dark patterns die u nog niet kende? Welke zijn het meest impactvol volgens u?}

Voornamelijk de systeemarchitect haalde enkele patronen aan waarvan hij het bestaan niet wist. Onder andere ``Toying with emotions'' en ``False hierarchy'' kwamen aan bod. Daarnaast beschouwen beide respondenten ``Obstruction'', ``Preselection'' en ``Hidden information'' als enorm impactvol omdat deze gebruikers kunnen aanzetten om ongewild een abonnement af te sluiten of langer dan gewenst een abonnement te behouden. Eveneens benadrukte de product designer dat het vooral de patronen zijn die op het eerste zicht niet schadelijk zijn die het meeste impact kunnen hebben, omdat gebruikers zich vaak niet bewust zijn van de invloed ervan.

\vspace{1em}

\textbf{4. Welke gevolgen verwacht u dat dark patterns op streamingdiensten hebben op gebruikers binnen de doelgroep van dit onderzoek?}

De respondenten verwachten zowel financiële als gedragsmatige gevolgen. Ze wijzen op het gevaar van ongewenste abonnementen, automatische ongewenste prijsverhogingen of upgrades en langer gebruik van streamingdiensten door functies zoals autoplay. Daarnaast kunnen patronen die inspelen op verslaving en het kijken op een mobieltoestel ervoor zorgen dat gebruikers op termijn veel minder buiten komen of praten met elkaar en een minder actieve levensstijl aannemen.

\newpage

\textbf{5. Welke soorten dark patterns van de acht meestvoorkomende op streamingdiensten vindt u het belangrijkst om te detecteren binnen een eerste proof of concept? Waarom?}

Voor een eerste proof of concept zou de product designer voornamelijk focussen op ``Preselection'', ``Hidden information'' en ``False hierarchy''. Volgens hem zijn dit patronen waarbij detectie en correcte informatieoverdracht een grote meerwaarde kunnen bieden, terwijl ze technisch haalbaar blijven binnen de scope van een proof of concept. Aanvullend beschouwen beide respondenten ``Roach motel'' als een belangrijk patroon om te detecteren. Echter erkennen ze dat dit technisch uitdagender is, waardoor het focussen op dit patroon in een latere fase wordt aangeraden. Ten slotte ziet de systeemarchitect ook een meerwaarde in het detecteren van ``Tricked Questions'' en de ``Countdown Timer''.

\vspace{1em}

\textbf{6. Welke uitdagingen verwacht u bij het automatisch detecteren van dark patterns?}

Volgens de product designer ligt de grootste uitdaging in het correct interpreteren van de context waarin een interface-element voorkomt. Niet elk opvallend element is namelijk automatisch een dark pattern, waardoor de onderliggende intentie moeilijk automatisch te bepalen is. Daarnaast ziet hij een uitdaging in het detecteren van dynamische patronen, aangezien hiervoor informatie over meerdere schermen of interactiemomenten heen moet worden bijgehouden. De systeemarchitect benadrukt aanvullend de juridische en beleidsmatige uitdagingen. Volgens hem is alles technisch haalbaar, maar zijn er beperkingen in wat wettelijk toegelaten is.

\vspace{1em}

\textbf{7. Op welke manier zou een tool gebruikers volgens u kunnen informeren over een mogelijk dark pattern?}

Beide respondenten vinden dat een tool gebruikers op een informatieve en niet-opdringerige manier moet ondersteunen. De nadruk wordt gelegd op een tool die zelf geen actie onderneemt maar die op de juiste momenten een korte, geargumenteerde uitleg geeft over patronen. Eveneens mag de tool de gebruikerservaring niet storen en moet er volgens de product designer nagedacht worden over de mogelijkheid tot resizing en verplaatsingsmogelijkheden. De systeemarchitect speelde eveneens met het idee van pas meldingen te geven buiten het venster waar content wordt afgespeeld. Ten slotte benadrukken beide respondenten dat gebruikers zelf moeten kunnen kiezen hoe en wanneer zij meldingen ontvangen.

\newpage

\textbf{8. Hoe zou u omgaan met situaties waarin de tool niet met zekerheid kan bepalen of een element een dark pattern is?}

De respondenten hadden verschillende standpunten over hoe omgegaan moet worden met onzekere detecties. Vanuit een technisch perspectief stelde de systeemarchitect voor om gebruik te maken van een betrouwbaarheidsniveau, bijvoorbeeld gebaseerd op een probability-score bij een Machine Learning-model. Wanneer deze score onder een bepaalde drempelwaarde ligt (bijvoorbeeld 70\%), zou de detectie volgens hem niet aan de gebruiker moeten worden weergegeven om foutieve waarschuwingen te vermijden. De product designer benadrukte daarentegen het belang van transparantie richting de gebruiker. Volgens hem is het beter om onzekere detecties wel te communiceren, maar duidelijk aan te geven dat het om een mogelijk dark pattern gaat. Daarnaast zouden gebruikers optioneel toegang moeten kunnen krijgen tot extra informatie over het patroon dat vermoed wordt om zelf een geïnformeerde keuze te maken. Ten slotte zou het krijgen van meldingen omtrent onzekere detecties ook optioneel moeten zijn.

\vspace{1em}

\textbf{9. Waar ligt volgens u de grens tussen informeren en storen?}

Beide respondenten geven aan dat de grens afhangt van de context en de mate van controle die de gebruiker heeft. Waarschuwingen mogen niet te frequent of onnodig verschijnen en moeten relevant blijven. De mogelijkheid om meldingen aan te passen of tijdelijk uit te schakelen wordt als een belangrijk aspect beschouwd.

\vspace{1em}

\textbf{10. Kan binnen deze use case in een beperkte mate het verstoren van de gebruiker gerechtvaardigd worden? Waarom wel of niet?}

De respondenten verschillen bij deze vraag van mening. De product designer vindt dat opvallendere waarschuwingen gerechtvaardigd kunnen worden indien de gebruiker een groot risico loopt, bijvoorbeeld bij het afsluiten van een abonnement met automatische verlenging. In dat geval zou de melding opvallender mogen ogen. Is het een patroon waarbij de gebruiker een minder groot risico loopt, dan toont men best een passieve melding. Anderzijds vindt de andere respondent dat de tool de gebruiker nooit actief mag storen en dat de voorkeuren van de gebruiker altijd gerespecteerd moeten worden.

\vspace{1em}

\textbf{11. Welke risico's ziet u wanneer gebruikers frequent waarschuwingen ontvangen?}

Beide respondenten verwachten dat een overlading aan waarschuwingen leidt tot gewenning, irritatie en het negeren van meldingen. In het slechtste geval zouden gebruikers de tool verwijderen. Daarom benadrukken zij dat waarschuwingen juist getimed, kwalitatief en configureerbaar moeten zijn.

\newpage

\textbf{12. Hoe kan volgens u een waarschuwing voldoende zichtbaar zijn zonder de gebruikerservaring te verstoren? Zou een gebruiker bijvoorbeeld controle moeten krijgen over de manier waarop waarschuwingen worden weergegeven? Waarom wel/niet?}

Beide respondenten vinden dat gebruikers voldoende controle moeten krijgen over de manier waarop waarschuwingen worden weergegeven. De systeemarchitect benadrukt dat gebruikers zelf moeten kunnen bepalen over welke dark patterns zij meldingen ontvangen en hoe frequent deze waarschuwingen worden weergegeven. De product designer legt daarnaast de nadruk op het belang van goede standaardinstellingen, waarbij enkel de noodzakelijke basisopties standaard geactiveerd zijn en gebruikers zelf verdere aanpassingen kunnen maken. Daarnaast raadt hij aan om te werken met subtiele en consistente meldingen in plaats van opvallende of schokkende waarschuwingen.

\vspace{1em}

\textbf{13. Welke eigenschappen zijn volgens u essentieel voor een succesvolle detectie- en signaleringstool?}

Beide respondenten benadrukken dat gebruikerscontrole een essentiële eigenschap is van een succesvolle detectie- en signaleringstool. Gebruikers moeten zelf kunnen bepalen hoeveel meldingen zij ontvangen, op welke momenten deze verschijnen en welke patronen zij willen opvolgen. Daarnaast benadrukt de product designer dat kwalitatieve detecties belangrijker zijn dan een grote hoeveelheid waarschuwingen. Volgens hem is het beter om minder detecties te tonen die betrouwbaar en relevant zijn dan gebruikers te overspoelen met meldingen. Verder wijst hij op het belang van transparantie, waarbij de tool een betrouwbaarheidsniveau, disclaimer en bijkomende uitleg voorziet bij een detectie.

\vspace{1em}

\textbf{14. Na het doorlopen van de huidige functionele en niet-functionele requirements, zijn er vereisten die volgens u ontbreken? Waar zou u al dan niet wijzigingen doorvoeren?}

Beide respondenten gaven enkele aandachtspunten mee met betrekking tot de huidige requirements. De opmerkingen kunnen als volgt worden samengevat:

\vspace{1em}

\textbf{S. Hostyn:}
\begin{itemize}
    \item Een onboarding ontbreekt momenteel als vereiste. Volgens hem is de eerste kennismaking met de tool belangrijk om vertrouwen op te bouwen en gebruikers te overtuigen van de meerwaarde ervan. Zonder onboarding kan er een drempel ontstaan om de tool effectief te gebruiken.
    \item Er zou een educatief aspect toegevoegd kunnen worden aan de requirements. Hierbij kan bijvoorbeeld gemeten worden of gebruikers door het gebruik van de tool effectief beter worden in het herkennen van dark patterns.
    \item De huidige scope met vier verschillende gebruikersrollen lijkt volgens hem te uitgebreid voor een proof of concept. De focus van de POC zou moeten liggen op het aantonen van de basisfunctionaliteiten en de interactie met de gebruiker. Het implementeren van alle rollen zou de omvang te sterk uitbreiden en eerder richting een volledig product evolueren.
\end{itemize}
    
\textbf{G. De Bruyn:}
\begin{itemize}
    \item De requirement waarbij een melding binnen één seconde na detectie van een dark pattern moet verschijnen, zou volgens hem weggehaald mogen worden. Voor een proof of concept is betrouwbaarheid en efficiëntie belangrijker dan een zeer korte responstijd.
    \item De requirement rond herstelbaarheid, waarbij minstens 95\% van de opgeslagen detectiegeschiedenis behouden moet blijven na een onverwachte applicatiecrash, vindt hij minder relevant binnen de scope van de bachelorproef en zou ook weggehaald mogen worden. 
\end{itemize}